\documentclass[11pt]{article}
\usepackage[a4paper,margin=25mm]{geometry}
\usepackage[T1]{fontenc}
\usepackage[utf8]{inputenc}
\usepackage{lmodern}
\usepackage{booktabs}
\usepackage{longtable}
\usepackage{array}
\usepackage{xspace}
\usepackage{hyperref}
\newcommand{\code}[1]{\texttt{#1}}
\newcommand{\us}{\textmu s\xspace}
\newcommand{\csitl}{CsI(Tl)\xspace}
\newcommand{\planned}{planned}
\newcommand{\doing}{in progress}
\newcommand{\done}{done}

\title{CryspLight full simulation: working plan}
\author{J.\,J.~Gómez-Cadenas}
\date{Live document --- last edited 16 October 2026}

\begin{document}
\maketitle

\section{Scope}

Build the full simulation specified in \code{fullsim\_design.tex} and compare
it with the bench measurements: BGO and \csitl at ambient temperature with the
Citiroc~1A (data taking the week of 20 October 2026), BGO at 196~K
afterwards, Radioroc2 when the board arrives. This document tracks what is
planned, in progress and done, the decisions taken and the questions open. It
is edited at every step; the physics stays in the design document.

\section{Principles}

\begin{itemize}
\item The pedagogical simulation is the git tag \code{pedsim} and the branch
  \code{main}; the full simulation grows on the branch \code{fullsim}.
  Nothing is replaced: every feature is gated by a configuration key and the
  frozen-output test proves the pedagogical path unchanged.
\item Repository layout for the full simulation: \code{docs\_fullsim/}
  (documents), \code{fullsim\_runs/<campaign>/} (simulation and analysis
  configurations plus the run log of each measurement campaign),
  \code{fullsim\_analysis/} (Python comparison layer, committed figures under
  \code{results/<campaign>/}), engine additions in \code{src/}, the
  electronics emulation in \code{src/frontend/}, tests in \code{test/}.
  Simulation output in the gitignored \code{output/<tag>/}. Measured data
  outside the repository, referenced by path from the campaign configuration.
\item Documents are written before the code they describe; results are
  produced by tracked scripts and committed with their images; every reported
  result states where it lives.
\end{itemize}

\section{Phases and tasks}

\subsection*{Phase 0 --- before and during the ambient runs}

\begin{longtable}{p{0.06\textwidth} p{0.62\textwidth} p{0.14\textwidth} p{0.1\textwidth}}
\toprule
id & task & status & date \\
\midrule
\endhead
0.1 & Tag \code{pedsim}, branch \code{fullsim}, datasheets committed & \done & 16 Oct \\
0.2 & Design document and this plan & \done & 16 Oct \\
0.3 & Overvoltage recipe for the bench (\code{BGOTSCAN/docs/overvoltage\_recipe.tex}) & \done & 16 Oct \\
0.4 & Frozen-output test of the pedagogical path & \planned & \\
0.5 & Readout layout: active rectangles, dead-area reflection, overhang; parity and layout tests & \planned & \\
0.6 & Grease and resin window: second Fresnel interface; window test & \planned & \\
0.7 & Materials: cold BGO file (2.7~\us), ambient BGO file; two-component \csitl timing & \planned & \\
0.8 & \Na point source with both lines and coincidence summing; source test & \planned & \\
0.9 & Photoelectron records \code{pes.h5}; re-binning test & \planned & \\
0.10 & SiPM device model (gain, PDE, crosstalk, afterpulsing, dark counts, pulse) as a module with unit tests & \planned & \\
0.11 & Citiroc chain: shapers, noise, trigger, 64 peak detectors, hold, ADC, saturation; unit tests & \planned & \\
0.12 & Common event layout writer (emulation) and the converter skeleton for measured files & \planned & \\
0.13 & First emulated campaign: BGO ambient and \csitl ambient at 2, 2.7, 3, 4~V overvoltage; tables of photoelectrons per channel, brightest-channel tail, stored fraction, resolution floor, trigger margin & \planned & \\
0.14 & Overvoltage scan figure and numbers delivered to the bench & \planned & \\
\bottomrule
\end{longtable}

\subsection*{Phase 1 --- when the ambient data exist}

\begin{longtable}{p{0.06\textwidth} p{0.62\textwidth} p{0.14\textwidth} p{0.1\textwidth}}
\toprule
id & task & status & date \\
\midrule
\endhead
1.1 & Converter from the bench files to the common layout, with the run log (bias, measured breakdown, thresholds, hold, shaping, temperatures) & \planned & \\
1.2 & Single-photoelectron gains and breakdown per channel; pedestals and noise term; dark rates & \planned & \\
1.3 & Spectra: 511, 1275, sum peak; resolution and its decomposition; 1275 to 511 ratio & \planned & \\
1.4 & Light yield in photoelectrons and in photons per MeV through the simulated collection & \planned & \\
1.5 & Maps: occupancy, mean charge, seed-channel photopeak; fit of the dead-area reflectivity and check of the roughness & \planned & \\
1.6 & Brightest-channel distribution and saturation in \csitl & \planned & \\
1.7 & Emulation rerun with the measured noise, gains and settings; side-by-side figures & \planned & \\
1.8 & Note section per campaign in the results document & \planned & \\
\bottomrule
\end{longtable}

\subsection*{Phase 2 --- accuracy and the next setups}

\begin{longtable}{p{0.06\textwidth} p{0.62\textwidth} p{0.14\textwidth} p{0.1\textwidth}}
\toprule
id & task & status & date \\
\midrule
\endhead
2.1 & Spectral sampling: emission spectra, $n(\lambda)$, PDE($\lambda$); spectral test & \planned & \\
2.2 & Cold BGO campaign at 196~K: breakdown measured in situ, emulation, comparison; warm-to-cold yield ratio against the quenching curve & \planned & \\
2.3 & Radioroc2 chain and its campaign & \planned & \\
2.4 & Gamma stage: K-shell fluorescence and escape & \planned & \\
2.5 & Position reconstruction on emulated maps: how much of the \code{pedsim} CNN resolution survives each front end & \planned & \\
2.6 & Design document version 1.0 with the measured parameters & \planned & \\
\bottomrule
\end{longtable}

\section{Campaigns}

\begin{tabular}{l l l l}
\toprule
campaign & configuration & readout & dates \\
\midrule
\code{2026-10\_citiroc\_ambient} & BGO ambient, \csitl ambient & Citiroc 1A $\times 2$ & week of 20 Oct 2026 \\
\code{2026-xx\_citiroc\_cold} & BGO 196~K & Citiroc 1A $\times 2$ & to be scheduled \\
\code{2026-xx\_radioroc} & all three & Radioroc2 & on delivery \\
\bottomrule
\end{tabular}

\section{Decisions}

\begin{description}
\item[16 Oct 2026] Full simulation as a subproject of CryspLight, not a client
  package: most of the work is engine work under the parity test; one history,
  one test suite.
\item[16 Oct 2026] Additive only. Tag \code{pedsim}; frozen-output test; gated
  features; separate files for the photoelectron records and the emulated
  events; pedagogical configurations and outputs untouched.
\item[16 Oct 2026] Three-stage pipeline with files between stages, so that the
  front end can be re-tuned against data without rerunning the optics.
\item[16 Oct 2026] Grease ($n = 1.42$) in every configuration; not a parameter.
\item[16 Oct 2026] Cold BGO decay time 2.7~\us from the Antarctica runs, in a
  new material file; the Ding file stays for the pedagogical runs.
\item[16 Oct 2026] The Antarctica runs are not simulated (dead channels, noisy
  chain); the first comparison is the Citiroc ambient campaign.
\item[16 Oct 2026] Overvoltage set from the measured breakdown, 3.0~V for BGO
  and 2.7~V for \csitl; recipe in \code{BGOTSCAN/docs/overvoltage\_recipe.tex}.
\item[16 Oct 2026] Measured and emulated events share one HDF5 layout; the
  bench files are converted to it as the first analysis step.
\item[16 Oct 2026] No waveform-digitiser emulation.
\end{description}

\section{Open questions}

\begin{itemize}
\item Reflectivity of the tile surface between channels (default diffuse 0.5;
  to fit on the occupancy maps).
\item BGO light yield at ambient and at 196~K for this crystal; BGO intrinsic
  resolution.
\item Citiroc settings of the ambient runs (shaping time, gains, hold delay,
  thresholds, trigger definition in the FPGA) and the ADC; to be written in
  the run log.
\item Angular dependence of the PDE through the resin from grease.
\item Citiroc shaper transfer function and noise spectrum; whether a
  single-photoelectron waveform can be probed on the test board.
\item Temperature of the SiPMs in the cold run (resolved by the measured
  breakdown; a thermometer on the SiPM board would still help).
\end{itemize}

\section{Log}

\begin{description}
\item[16 Oct 2026] Project assessed; datasheets read (Citiroc 1A, Radioroc2,
  S14160/S14161); Antarctica draft read; tag and branch; design document
  v0.1 and this plan; overvoltage recipe written.
\end{description}

\end{document}
